\documentclass[10pt, aspectratio=169]{beamer}
\usepackage{../beamer_theme_408}

\title{408考研零基础 --- 数据结构}
\subtitle{1-4 线性表的链式存储}
\author{}
\date{}


\begin{document}


\begin{frame}{本节目标}
    \begin{enumerate}
        \item 理解单链表的结点结构，掌握头指针与头结点的概念
        \item 掌握单链表的建立方法（头插法和尾插法）
        \item 掌握单链表的插入、删除、查找操作的实现与时间复杂度分析
    \end{enumerate}
\end{frame}

\begin{frame}[fragile]{单链表的结构定义}
    每个结点包含数据域和指针域
    \vspace{0.5em}
    \begin{lstlisting}[language=C, basicstyle=\ttfamily\small]
typedef struct LNode {
    ElemType data;
    struct LNode *next;
} LNode, *LinkList;
    \end{lstlisting}
    \vspace{0.5em}
    \textbf{重要概念}
    \begin{itemize}
        \item \textbf{头指针}：指向链表第一个结点的指针
        \item \textbf{头结点}：第一个结点前的附加结点
    \end{itemize}
    \vspace{0.3em}
    \begin{block}{引入头结点的好处}
        第一个元素前插入或删除第一个元素时，操作与其他位置一致
    \end{block}
\end{frame}

\begin{frame}{单链表的建立 --- 头插法}
    \textbf{每次将新结点插入到头结点之后}
    \vspace{0.5em}
    \begin{enumerate}
        \item 生成新结点 $s$
        \item $s \to \text{next} = \text{head} \to \text{next}$
        \item $\text{head} \to \text{next} = s$
    \end{enumerate}
    \vspace{0.5em}
    \begin{block}{特点}
        链表顺序与输入顺序相反（逆序建立）
    \end{block}
\end{frame}

\begin{frame}{单链表的建立 --- 尾插法}
    \textbf{每次将新结点插入到链表尾部}
    \vspace{0.5em}
    维护尾指针 $r$ 指向最后一个结点
    \vspace{0.5em}
    \begin{enumerate}
        \item 生成新结点 $s$
        \item $r \to \text{next} = s$
        \item $r = s$
    \end{enumerate}
    \vspace{0.5em}
    \begin{block}{特点}
        链表顺序与输入顺序相同（正序建立）
    \end{block}
\end{frame}

\begin{frame}{单链表的查找}
    \textbf{按位查找}
    \begin{itemize}
        \item 从头结点开始，沿 next 指针移动
        \item 时间复杂度 $O(n)$
    \end{itemize}
    \vspace{0.5em}
    \textbf{按值查找}
    \begin{itemize}
        \item 从第一个元素结点开始依次比较
        \item 时间复杂度 $O(n)$
    \end{itemize}
    \vspace{0.5em}
    \begin{alertblock}{与顺序表对比}
        顺序表按位查找 $O(1)$，链表 $O(n)$
    \end{alertblock}
\end{frame}

\begin{frame}{单链表的插入操作}
    \textbf{后插（已知结点 $p$）}
    \[
    s \to \text{next} = p \to \text{next};\quad p \to \text{next} = s;
    \]
    时间复杂度 $O(1)$
    \vspace{0.5em}
    \textbf{前插（已知结点 $p$）}
    \begin{itemize}
        \item 后插 + 交换数据
        \item 时间复杂度 $O(1)$
    \end{itemize}
    \vspace{0.5em}
    \textbf{按位序插入}
    \begin{itemize}
        \item 先查找第 $i-1$ 个结点
        \item 再后插
        \item 时间复杂度 $O(n)$
    \end{itemize}
\end{frame}

\begin{frame}{单链表的删除操作}
    \textbf{删除 $p$ 的后继 $q$}
    \[
    p \to \text{next} = q \to \text{next};\quad \text{free}(q);
    \]
    时间复杂度 $O(1)$
    \vspace{0.5em}
    \textbf{删除指定结点 $p$}
    \begin{itemize}
        \item 将后继数据复制到 $p$，再删除后继
        \item 时间复杂度 $O(1)$
        \item 注意：$p$ 为尾结点时无效
    \end{itemize}
    \vspace{0.5em}
    \textbf{按位序删除}
    \begin{itemize}
        \item 先查找第 $i-1$ 个结点
        \item 时间复杂度 $O(n)$
    \end{itemize}
\end{frame}

\begin{frame}{例题 --- 单链表就地逆置}
    \textbf{设计算法，对带头结点的单链表就地逆置}
    \[
    O(n), O(1)
    \]
    \vspace{0.3em}
    \textbf{思路}：头插法
    \begin{enumerate}
        \item $\text{head} \to \text{next} = \text{NULL}$
        \item 遍历原链表，依次取下每个结点 $p$
        \item 使用头插法将 $p$ 插入 head 之后
    \end{enumerate}
\end{frame}

\begin{frame}{本节总结}
    \begin{enumerate}
        \item 单链表结点由数据域和指针域组成
        \item 头指针标识链表，头结点简化边界处理
        \item 建立方式：头插法（逆序）、尾插法（正序）
        \item 查找 $O(n)$，已知位置插入删除 $O(1)$
        \item 链表适用于插入删除频繁而查找较少的场景
    \end{enumerate}
\end{frame}

\end{document}
